% ---------- Datei: chapters/03_konzeption.tex ----------
\chapter{Konzeptionsphase}

\section{Technologieauswahl}
Zur Erfüllung der Anforderungen werden folgende Google Cloud Technologien eingesetzt:
\begin{itemize}
    \item Network Connectivity Center (NCC)
    \item VPC Networks und Cloud Router
    \item Cloud VPN oder Cloud Interconnect
    \item Cloud Firewall und Cloud Monitoring
\end{itemize}

\section{Netzwerkdesign und NCC-Architektur}
Der NCC Hub bildet das Kernstück der neuen Architektur und fungiert als zentrale Verwaltungsplattform. Zwei VPC-Netzwerke in unterschiedlichen Regionen sowie zwei Zweigstellen werden als sogenannte Spokes an den Hub angebunden. Dies ermöglicht eine automatische Routenverteilung zwischen allen Spokes via Border Gateway Protocol (BGP).

\section{Sicherheitskonzept und IP-Adressierung}
Um überlappende Subnetze zwischen der Cloud und den lokalen Standorten zu vermeiden, wird ein strukturierter IP-Adressplan entworfen. Zusätzlich werden Firewall-Richtlinien konzipiert, die festlegen, welche Netzsegmente untereinander kommunizieren dürfen.